\documentclass[../numerical_methods.tex]{subfiles}
\begin{document}
\section{Aula 16 - 02 de Junho, 2026}
\subsection{Motivações}
\begin{itemize}
	\item Método de Francis
\end{itemize}
\subsection{Método de Francis}
Seja \(A\in \mathbb{M}_{n}\) uma matriz simétrica; vamos usar a decomposição QR para calcular todos os seus autovalores e autovetores:
\begin{align*}
	 & A_1 = A \Rightarrow A_1 = Q_1 R_1                             \\
	 & A_2 = R_1Q_1 \Rightarrow A_2 = Q_2R_2                         \\
	 & A_3 = R_2Q_2 \Rightarrow A_3 = Q_3R_3                         \\
	 & \quad \vdots                                                  \\
	 & A_{k-1} = R_{k-2}Q_{k-2} \Rightarrow A_{k-1} = Q_{k-1}R_{k-1} \\
	 & A_{k} = R_{k-1}Q_{k-1}.
\end{align*}
Mostraremos duas coisas:
\begin{itemize}
	\item[i)] Todas as matrizes desta sequência são semelhantes; e
	\item[ii)] As matrizes\(A_{k}\) convergem para uma matriz diagonal.
\end{itemize}
Começando do final,
\begin{prop*}
	As matrizes A e \(A_{k}\) possuem os mesmos autovalores; além disso, os autovetores estarão contidos no produto de todas as \(Q_{i},\; 1\leq i\leq k-1\).
\end{prop*}
\begin{proof*}
	Como \(A_{k-1}=Q_{k-1}R_{k-1}\), vale também que
	\[
		R_{k-1} = Q_{k-1}^{T}A_{k-1} \quad\&\quad A_{k} = R_{k-1}Q_{k-1}.
	\]
	Daí, substituindo em \(R_{k-1}\), temos \(A_{k} = Q_{k-1}^{T}A_{k-1}Q_{k-1}\) e, consequentemente,
	\begin{align*}
		A_{k} & = Q_{k-1}^{T}A_{k-1}Q_{k-1}                                          \\
		      & = Q_{k-1}^{T}Q_{k-2}^{T}A_{k-2}Q_{k-2}Q_{k-1}                        \\
		      & \vdots                                                               \\
		      & = Q_{k-1}^{T}Q_{k-2}^{T}\dotsc Q_{1}^{T}A_1Q_1 \dotsc Q_{k-2}Q_{k-1} \\
		      & = (Q_1 \dotsc Q_{k-2}Q_{k-1})^{T}A (Q_1 \dotsc Q_{k-2}Q_{k-1}).
	\end{align*}
	Assim, fazendo \(V=Q_1 \dotsc Q_{k-2}Q_{k-1}\), obtemos a relação de A e \(A_{k} = V^{T}AV\) serem semelhantes. Portanto, os autovalores de A e \(A_{k}\) são os mesmos, com os autovetores contidos na matriz V. \qedsymbol
\end{proof*}
\begin{prop*}
	A sequência \(A_{k}\) converge para uma matriz diagonal.
\end{prop*}
Logo, os elementos da diagonal de \(A_{k}\) fornecem uma aproximação para os autovalores de A, enquanto as colunas da matriz \(V = Q_1Q_2 \dotsc Q_{k-1}\) fornecem uma aproximação dos seus respectivos autovetores.
Implementando em código,
\begin{listing}[H]
	\caption{Implementação do Método de Francis e Comparações}
	\begin{minted}[bgcolor = DarkBackground, linenos=true, breaklines]{Python}
import numpy as np

# QR Clássico
def clgsQR(A):
    (m, n) = np.shape(A)
    Q = np.zeros((m, n))
    R = np.zeros((n, n))

    for j in range(n):
        V = A[:, j].copy()
        for i in range(j):
            R[i, j] = Q[:, i].dot(A[:, j])
            V = V - R[i, j] * Q[:, i]
        R[j, j] = np.linalg.norm(V)
        Q[:, j] = V / R[j, j]

    return Q, R
 \end{minted}
\end{listing}
\begin{listing}[H]
	\begin{minted}[bgcolor = DarkBackground, linenos=true, breaklines]{Python}
# QR Modificado
def mgsQR(A):
    (m, n) = np.shape(A)
    V = np.copy(A).astype(float) # Força float para evitar problemas de tipo
    Q = np.zeros((m, n))
    R = np.zeros((n, n))

    for j in range(n):
        for i in range(j):
            R[i, j] = Q[:, i].dot(V[:, j])
            V[:, j] = V[:, j] - R[i, j] * Q[:, i]
        R[j, j] = np.linalg.norm(V[:, j])
        Q[:, j] = V[:, j] / R[j, j]

    return Q, R
 \end{minted}
\end{listing}
\begin{listing}[H]
	\begin{minted}[bgcolor = DarkBackground, linenos=true, breaklines]{Python}
def francis(A, tol, flag):
    if flag == 'classico':
        print('QR', flag, '\n')
        qr = clgsQR
    elif flag == 'modificado':
        print('QR', flag, '\n')
        qr = mgsQR
    elif flag == 'python':
        print('QR', flag, '\n')
        qr = np.linalg.qr

    n = np.shape(A)[0]
    A_local = np.copy(A).astype(float) # Garante que estamos tratando como float
    V = np.eye(n)
    erro = np.inf
 \end{minted}
\end{listing}
\begin{listing}[H]
	\begin{minted}[bgcolor = DarkBackground, linenos=true, breaklines]{Python}
    # Trava de segurança para não travar o computador em loops infinitos
    max_iter = 5000
    it_atual = 0

    while erro > tol and it_atual < max_iter:
        [Q, R] = qr(A_local)

        A_local = R.dot(Q)
        V = V.dot(Q)

        # Calcula o erro olhando o maior valor absoluto abaixo da diagonal principal
        erro = np.max(np.abs(np.tril(A_local, -1)))
        it_atual += 1

    if it_atual == max_iter:
        print(f"Aviso: O método atingiu o limite de {max_iter} iterações sem convergir totalmente.")
    else:
        print(f"Convergiu em {it_atual} iterações.")

    D = np.diag(A_local)
    return D, V

 \end{minted}
\end{listing}
Para testar o código, temos alguns exemplos.

\begin{listing}[H]
	\begin{minted}[bgcolor = DarkBackground, linenos=true, breaklines]{Python}
# ========================================================
# EXEMPLO INTERESSANTE: Construindo a Matriz de Teste
# ========================================================
epsilon = 1e-4
B = np.array([[1,       1,       1      ],
              [epsilon, 0,       0      ],
              [0,       epsilon, 0      ],
              [0,       0,       epsilon]], dtype=float)

# Multiplicamos B^T * B para gerar uma matriz simétrica para o método de Francis

C = np.transpose(B).dot(B)

print("--- Matriz de Teste C (Simétrica e Mal-Condicionada) ---")
print(C)
print("-" * 60)

tol = 1e-5

\end{minted}
	\begin{minted}[bgcolor = DarkBackground, linenos=true, breaklines]{Python}
# ------------------------------------------------------------
# Teste 1: Francis com Gram-Schmidt Clássico
# ------------------------------------------------------------
print("\n>>> Executando Francis com Gram-Schmidt CLÁSSICO...")
D_c, V_c = francis(C, tol, flag='classico')

print("Autovalores encontrados (CGS):", D_c)
I_c = np.transpose(V_c).dot(V_c)
erro_ort_c = np.linalg.norm(np.eye(len(C)) - I_c, 'fro')
print(f"Erro de Ortogonalidade dos Autovetores (V): {erro_ort_c:.2e}")
print("-" * 60)
\end{minted}
\end{listing}
\begin{listing}[H]
	\begin{minted}[bgcolor = DarkBackground, linenos=true, breaklines]{Python}
# ------------------------------------------------------------
# Teste 2: Francis com Gram-Schmidt Modificado
# ------------------------------------------------------------
print("\n>>> Executando Francis com Gram-Schmidt MODIFICADO...")
D_m, V_m = francis(C, tol, flag='modificado')

print("Autovalores encontrados (MGS):", D_m)
I_m = np.transpose(V_m).dot(V_m)
erro_ort_m = np.linalg.norm(np.eye(len(C)) - I_m, 'fro')
print(f"Erro de Ortogonalidade dos Autovetores (V): {erro_ort_m:.2e}")
print("-" * 60)
\end{minted}
	\begin{minted}[bgcolor = DarkBackground, linenos=true, breaklines]{Python}
# ------------------------------------------------------------
# Teste 3: Francis com o QR nativo do Python (NumPy)
# ------------------------------------------------------------
print("\n>>> Executando Francis com o QR nativo do PYTHON...")
D_p, V_p = francis(C, tol, flag='python')

print("Autovalores encontrados (Python):", D_p)
I_p = np.transpose(V_p).dot(V_p)
erro_ort_p = np.linalg.norm(np.eye(len(C)) - I_p, 'fro')
print(f"Erro de Ortogonalidade dos Autovetores (V): {erro_ort_p:.2e}")
print("-" * 60)
\end{minted}
\end{listing}

Agora, para finalizarmos, um exemplo com matriz cheia:
\begin{listing}[H]
	\begin{minted}[bgcolor = DarkBackground, linenos=true, breaklines]{Python}
# =========================================
# NOVO EXEMPLO: Uma matriz cheia
# =========================================
# Esta matriz não é diagonal e vai exigir que o QR esmague os elementos inferiores
C = np.array([[6.0, 2.0, 1.0],
              [2.0, 5.0, 2.0],
              [1.0, 2.0, 7.0]], dtype=float)

print("--- Nova Matriz de Teste C (Cheia e Densa) ---")
print(C)
print("-" * 60)

tol = 1e-12

\end{minted}
	\begin{minted}[bgcolor = DarkBackground, linenos=true, breaklines]{Python}
# -------------------------------------------------
# Teste 1: Francis com Gram-Schmidt Clássico
# -------------------------------------------------
print("\n>>> Executando Francis com Gram-Schmidt CLÁSSICO...")
D_c, V_c = francis(C, tol, flag='classico')

print("Autovalores encontrados (CGS):\n", D_c)
I_c = np.transpose(V_c).dot(V_c)
erro_ort_c = np.linalg.norm(np.eye(len(C)) - I_c, 'fro')
print(f"Erro de Ortogonalidade dos Autovetores (V): {erro_ort_c:.2e}")
print("-" * 60)
\end{minted}
\end{listing}
\begin{listing}[H]
	\begin{minted}[bgcolor = DarkBackground, linenos=true, breaklines]{Python}
# ---------------------------------------------------
# Teste 2: Francis com Gram-Schmidt Modificado
# ---------------------------------------------------
print("\n>>> Executando Francis com Gram-Schmidt MODIFICADO...")
D_m, V_m = francis(C, tol, flag='modificado')

print("Autovalores encontrados (MGS):\n", D_m)
I_m = np.transpose(V_m).dot(V_m)
erro_ort_m = np.linalg.norm(np.eye(len(C)) - I_m, 'fro')
print(f"Erro de Ortogonalidade dos Autovetores (V): {erro_ort_m:.2e}")
print("-" * 60)
\end{minted}
	\begin{minted}[bgcolor = DarkBackground, linenos=true, breaklines]{Python}
# --------------------------------------------------
# Teste 3: Francis com o QR nativo do Python (NumPy)
# --------------------------------------------------
print("\n>>> Executando Francis com o QR nativo do PYTHON...")
D_p, V_p = francis(C, tol, flag='python')

print("Autovalores encontrados (Python):\n", D_p)
I_p = np.transpose(V_p).dot(V_p)
erro_ort_p = np.linalg.norm(np.eye(len(C)) - I_p, 'fro')
print(f"Erro de Ortogonalidade dos Autovetores (V): {erro_ort_p:.2e}")
print("-" * 60)
\end{minted}
\end{listing}
\end{document}

